% NCSSL positioning — reusable paragraph for the intro / related-work handoff.
% Rests only on the TASK axis (imputation vs. downstream classification), so it
% needs no assumption about the previous paper's datasets or feature types.
% Framing rule: extends/scopes, never "overturns" or "redundant" near [NCSSL].

Our prior work~[NCSSL] introduced neighborhood-centric self-supervision for graph
attribute \emph{imputation}, where predicting neighborhood structure directly aids
the recovery of missing feature values. In this paper we ask a distinct question:
does that auxiliary signal also improve the \emph{downstream classification}
boundary, and under what conditions? These are different objectives---faithful
feature reconstruction is neither necessary nor sufficient for an accurate
decision boundary---so the answer is not implied by the earlier result. We find
that the transfer of neighborhood self-supervision from reconstruction to
downstream utility is \textbf{regime-dependent}: it is most valuable on
\textbf{dense-embedding} graphs and under \textbf{scarce} observations, while on
sparse binary attributes the neighborhood targets that aid reconstruction carry
little class-discriminative signal. This delineates \emph{where} neighborhood
self-supervision transfers to the downstream task, extending---rather than
overturning---the earlier framework.
